\section{谱序列的构造}
\label{sec:construction-spectral-sequence}

\stfhIndexRangeStart{I179}{Puxulie}{谱序列（spectral sequence）}
本节介绍一个重要工具——谱序列。它对计算\personFloer{}同调或把\personFloer{}同调用于实际问题
都很有用。这个特定谱序列把拉格朗日子流形 $L$ 的内蕴拓扑与嵌入
$L\subset(M,\omega)$ 的外在辛拓扑联系起来；作者在 (Oh96b) 中用小
$\epsilon>0$ 时的\personFloer{}复形
\[
 CF(\phi_{-\epsilon f}^1(L),L)
\]
（取 $\mathbb Z_2$ 系数）表述了它。\personBuhovsky{} (Bu10)、\personBiran{}--\personOctavianCornea{} (BCo09)
和\personDamian{} (D09) 仍以小\personMorse{}函数的\personMorse{}同调为基础，加强了这一构造。不过，
这个谱序列的最佳设定，是 (FOOO09) 通过直接考察干净对子 $(L,L)$ 而给出的\personBott{}--\personMorse{}设定。后一构造所需数学超出本书框架，故我们只讨论 (Oh96b) 的设定；
但至少会说明一个重要成分，即谱序列乘积结构的相容性。这里借用 (FOOO09) 定理 D
的叙述，并作适合单调情形的简化。

由于\personFloer{}同调中要用到乘积结构，本节采用\personFloer{}同调的上同调版本；这意味着把
$-\epsilon f$ 换成 $\epsilon f$。

基于\termThickThinDichotomy{}和局部\personFloer{}上同调的分析，固定\personMorse{}函数
$f:L\to\mathbb R$ 和充分小的 $\epsilon>0$，使得命题 17.1.3 与定理 17.2.7 的
二分对 $(L,\phi_{\epsilon f}^1(L))$ 成立。以 $\Sigma_L\geq2$ 表示 $L$ 的最小
正马斯洛夫数。

置 $N_L=[(n+1)/\Sigma_L]$，并如前定义形式参数
\[
 q^\beta=T^{\omega(\beta)}e^{\mu_L(\beta)},
 \qquad
 \deg(q^\beta)=\mu_L(\beta),\quad E(q^\beta)=\omega(\beta),
\]
再考虑\personNovikov{}环 $\Lambda_{(\omega,L)}$ 及其子环
$\Lambda_{0,(\omega,L)}$。

考虑\personFloer{}复形 $CF(L,\phi_{\epsilon f}^1(L))$，并按第 13.9.3 节的方式赋予它
$\Lambda_{(L,\omega)}$-模结构。回忆，第 13.9.3 节把与完整道路空间
$\Omega(L_0,L_1)$ 相伴的
$CF^*(L_1,L_0)\cong CF_*(L_0,L_1)^*$ 看成 $\Lambda_{(L_0,L_1)}$-模。

在当前 $(L,\phi_{\epsilon f}^1(L))$ 情形中，把 $q^\beta[p]$ 的滤过层级与分次定义为
\begin{equation}
 \ell(q^\beta[p])=\omega(\beta)+f(p),
\label{eq:17-filtration-q-beta-p}
\end{equation}
\begin{equation}
 \deg(q^\beta[p])=n-\mu_f(p)-\mu_L(\beta).
\label{eq:17-grading-q-beta-p}
\end{equation}
这样设定，是为了让 $\overline\delta_{(0)}$ 与 $-f$ 的\personMorse{}边界算子（或 $f$ 的\personMorse{}上边界算子）一致。换言之，在 $\Lambda_{(\omega,L)}$ 中延拓关系
\begin{equation}
 [\beta]\cdot[p]=T^{-\omega(\beta)}e^{-\mu_L(\beta)}\cdot[p],
\label{eq:17-module-relation-beta-p}
\end{equation}
由此给 $CF(L,\phi_{\epsilon f}^1(L))$ 赋予 $\Lambda_{(\omega,L)}$-模结构。按定义，
\[
 CF(L,\phi_{\epsilon f}^1(L))=\Lambda_{\omega,L}\{\operatorname{Crit}f\}.
\]
于是上边界映射
\[
 \delta:CF(L,\phi_{\epsilon f}^1(L))\longrightarrow
 CF(L,\phi_{\epsilon f}^1(L))
\]
保持这个 $\Lambda_{(\omega,L)}$-模结构。

由\termThickThinDichotomy{}，边界算子 $\delta:C^*\to C^*$ 具有形式
\[
 \delta=\delta_{(0)}+\delta',
\]
其中 $\delta_{(0)}$ 来自与局部\personFloer{}复形相伴的上边界算子，$\delta'$ 是厚流线的
贡献。

为进一步分析 $\delta'$，设
$u\in\mathcal M(L,\phi_{\epsilon f}^1(L);x,y)$ 是任意厚流线，
$w:([0,1]^2,\partial[0,1]^2)\to(M,L)$ 是引理 17.1.4 证明中使用的相伴映射。
若 $\epsilon>0$ 取得充分小，则
\[
 \int w^*\omega\geq\Sigma_L-\frac14\Sigma_L
 =\frac34\Sigma_L>0,
 \qquad
 \mu_L(w)=1-\mu_f(x)+\mu_f(y)\leq n+1.
\]
\footnote{译者注：原书在已把 $\Sigma_L$ 定义为最小马斯洛夫数后，直接在面积
不等式中用它与 $\int w^*\omega$ 比较。除非另有未声明的单调常数归一化，否则两边
量纲不合；因所缺比例常数未在此固定，译文照录而不作不可唯一确定的修补。}
由 $L$ 的单调性与 $\int w^*\omega>0$，有 $\mu_L(w)>0$。因而
\[
 \mu_f(x)-\mu_f(y)=1-\mu_L(w),
\]
且
\[
 \mu_L(w)=l\Sigma_L,
 \qquad 1\leq l\leq\left[\frac{n+1}{\Sigma_L}\right].
\]
所以，每条厚流线都把\personMorse{}分次改变 $-l\Sigma_L+1$，其中
$1\leq l\leq[(n+1)/\Sigma_L]$；于是 $\delta$ 分解为
\begin{equation}
 \delta'=\overline\delta_{(1)}\otimes q
 +\cdots+\overline\delta_{(N_L)}\otimes q^{N_L},
\label{eq:17-delta-prime-decomposition}
\end{equation}
这里 $\overline\delta_{(l)}:C^{(*)}(f)\to C^{(*-l\Sigma_L+1)}(f)$ 是由连接
指标为 $*$ 与 $*-l\Sigma_L+1$ 的 $f$ 的临界点的流线所诱导的线性映射。回忆
$\delta_{(0)}$ 把 $C^*$ 映到 $C^{*+1}$。
\footnote{译者注：原书把目标次数印作 $*-1+l\Sigma_L$，与上一句给出的分次变化
$-l\Sigma_L+1$ 以及 $l=0$ 时 $\delta_{(0)}:C^*\to C^{*+1}$ 均矛盾；译文据这
两处改为 $*-l\Sigma_L+1$。}
\footnote{译者注：原书先把 $\delta'$ 定义为厚流线的贡献，但公式 (17.3.14) 的
右端又从局部项 $\overline\delta_{(0)}$ 开始；随后按 $q$ 次数展开的论证唯一确定
厚项应从 $l=1$ 开始，译文据此删去该局部项。}

\begin{theorem}
\label{thm:lagrangian-spectral-sequence}
存在具有如下性质的谱序列：
\begin{enumerate}[label=(\arabic*)]
\item
\[
 E_2^{p,q}=\bigoplus_k H^k(L;\mathbb Q)\otimes
 \left(T^{q\lambda}\Lambda_{(\omega,L)}
 /T^{(q+1)\lambda}\Lambda_{(\omega,L)}\right)^{(p-k)},
\]
其中 $\lambda>0$；
\item \personFloer{}上同调 $HF(L,\phi_{\epsilon f}^1(L);\Lambda_{(\omega,L)})$ 上存在滤过
$F^*HF(L,\phi_{\epsilon f}^1(L);\Lambda_{(\omega,L)})$，使得
\[
 E_\infty^{p,q}\cong
 \frac{F^qHF^p(L,\phi_{\epsilon f}^1(L);\Lambda_{(\omega,L)})}
 {F^{q+1}HF^p(L,\phi_{\epsilon f}^1(L);\Lambda_{(\omega,L)})}.
\]
\end{enumerate}
\end{theorem}

（更一般的陈述见 (FOOO09) 定理 D，尤其是由 \emph{bulk 形变}或由 $M$ 的环境上同调类
形变的版本。）本小节余下部分证明这个定理。

因为 $\Sigma_L\geq2$，$\overline\delta_{(l)}$ 的次数不大于 $-1$。比较方程
\[
 0=\delta\circ\delta=(\delta_{(0)}+\delta')^2
\]
各项的次数，立即得到
\[
 \delta_{(0)}\circ\delta_{(0)}=0,
 \qquad
 \delta_{(0)}\circ\delta'+\delta'\circ\delta_{(0)}
 +\delta'\circ\delta'=0.
\]
第一式使我们可以定义
$H(C_*,\delta_{(0)})\cong H^*(L;f)$。利用
\eqref{eq:17-delta-prime-decomposition} 展开第二式：
\[
 \begin{aligned}
 0&=\delta_{(0)}\circ\delta'+\delta'\circ\delta_{(0)}
 +\delta'\circ\delta'\\
 &=\bigl(\delta_{(0)}\overline\delta_{(1)}
 +\overline\delta_{(1)}\delta_{(0)}\bigr)q\\
 &\quad+\bigl(\overline\delta_{(1)}\overline\delta_{(1)}
 +\delta_{(0)}\overline\delta_{(2)}
 +\overline\delta_{(2)}\delta_{(0)}\bigr)q^2+\cdots.
 \end{aligned}
\]
收集 $q$-次数为 1 和 2 的项，得到
\[
 \delta_{(0)}\overline\delta_{(1)}+
 \overline\delta_{(1)}\delta_{(0)}=0,
\]
\[
 \overline\delta_{(1)}\overline\delta_{(1)}
 +\delta_{(0)}\overline\delta_{(2)}
 +\overline\delta_{(2)}\delta_{(0)}=0.
\]
第一式说明 $\overline\delta_{(1)}$ 下降到
$H(C_*,\delta_{(0)})=H^*(L;\epsilon f)$ 上，第二式说明它在那里定义一个微分。

$C(L,\phi_{\epsilon f}^1(L))$ 上的作用量滤过诱导出所定义的谱序列。首先有同构
\[
 \operatorname{gr}_*CF(L,\phi_{\epsilon f}^1(L))
 \cong C(L;\epsilon f)\otimes
 \operatorname{gr}_*(F^*\Lambda_{(\omega,L)})
\]
作为 $\operatorname{gr}_*(F^*\Lambda_{(\omega,L)})$-模。（关于
$\Lambda_{(\omega,L)}$ 上的滤过，见定义 16.2.3 及其后。）于是得到谱序列
$E_r^{p,q}$，其中
\begin{equation}
 E_2^{p,q}=\bigoplus_{k=0}^n H^k(L;R)\otimes
 \operatorname{gr}_q\bigl(F\Lambda_{(\omega,L)}^{(p-k)}\bigr).
\label{eq:17-spectral-sequence-E2}
\end{equation}
这里 $\Lambda_{(\omega,L)}^{(p-k)}$ 表示次数 $p-k$ 的部分。还可证明
$HF^p(L,\phi_{\epsilon f}^1(L);\Lambda_{(\omega,L)})$ 上存在滤过，使得
\[
 E_\infty^{p,q}\cong
 F^qHF^p(L,\phi_{\epsilon f}^1(L);\Lambda_{(\omega,L)})
 /F^{q+1}HF^p(L,\phi_{\epsilon f}^1(L);\Lambda_{(\omega,L)}).
\]
上述谱序列的收敛 $E_2^{p,q}\Longrightarrow E_\infty^{p,q}$ 的证明对任意交换环
$R$ 都成立。

现在考虑带滤过 $F^*D$ 的分次 $\Lambda_{(\omega,L)}$-模 $D$。该滤过与
$\Lambda_{(\omega,L)}$ 上的滤过相容；对
$D=CF(L,\phi_{\epsilon f}^1(L))$，即
\[
 F^l\Lambda_{(\omega,L)}\cdot F^kD\subseteq F^{k+l}D.
\]
令 $\Lambda_{(\omega,L)}^{(0)}$ 为 $\Lambda_{(\omega,L)}$ 的零次部分。则
\[
 \Lambda_{(\omega,L)}^{(0)}
 \cong R[[T^{A(\omega,L)}]][T^{-A(\omega,L)}].
\]
这说明，若 $R$ 是域，则 $\Lambda_{(\omega,L)}^{(0)}$ 也是域，因此 $D$ 的 $p$ 次
部分 $D^p$ 是 $\Lambda_{(\omega,L)}^{(0)}$ 上的向量空间。依照滤过复形的标准过程
(McC85)，用这个滤过定义谱序列。

注意，分别乘以 $T^{A(\omega,L)}$ 与 $e$ 给出同构
\[
 T^{A(\omega,L)}\cdot:F^qHF^p(L,\phi_{\epsilon f}^1(L);\Lambda_{(\omega,L)})
 \longrightarrow F^{q+1}HF^p(L,\phi_{\epsilon f}^1(L);\Lambda_{(\omega,L)}),
\]
\[
 e\cdot:F^qHF^p(L,\phi_{\epsilon f}^1(L);\Lambda_{(\omega,L)})
 \longrightarrow F^qHF^{p+2}(L,\phi_{\epsilon f}^1(L);\Lambda_{(\omega,L)}).
\]
记
\[
 \Lambda^{(0)}(\lambda)=\Lambda_{(\omega,L)}^{(0)}
 /F^\lambda\Lambda_{(\omega,L)}^{(0)}.
\]
以
$F^n\Lambda_{(\omega,L)}^{(0)}=F^{nA(\omega,L)}\Lambda_{(\omega,L)}^{(0)}$
定义 $\Lambda_{(\omega,L)}^{(0)}$ 的滤过。其相伴分次模为
\[
 \operatorname{gr}_*\bigl(F\Lambda_{(\omega,L)}^{(0)}\bigr)
 =\bigoplus_{n\in\mathbb Z_{\geq0}}
 \operatorname{gr}_n\bigl(F\Lambda_{(\omega,L)}^{(0)}\bigr),
\]
其中每个 $\operatorname{gr}_n(F\Lambda_{(\omega,L)}^{(0)})$ 自然同构于
$\Lambda^{(0)}(A(\omega,L))$。又有
\[
 \operatorname{gr}_*(F\Lambda_{(\omega,L)})
 =\operatorname{gr}_*(F\Lambda_{(\omega,L)}^{(0)})[e,e^{-1}]
 =\bigoplus_{n\in\mathbb Z_{\geq0}}
 \operatorname{gr}_n(F\Lambda_{(\omega,L)}^{(0)})[e,e^{-1}].
\]

以 $\widehat C$ 表示 $\operatorname{Crit}f$ 上的
$\Lambda_{(\omega,L)}^{(0)}$-模，并定义
\[
 \begin{aligned}
 Z_r^{p,q}(\widehat C)
 &=\{x\in F^q\widehat C^p\mid
 \delta x\in F^{q+r-1}\widehat C^{p+1}\}+F^{q+1}\widehat C^p,\\
 B_r^{p,q}(\widehat C)
 &=\bigl(\delta(F^{q-r+2}\widehat C^{p-1})\cap F^q\widehat C^p\bigr)
 +F^{q+1}\widehat C^p,\\
 E_r^{p,q}(\widehat C)&=Z_r^{p,q}(\widehat C)/B_r^{p,q}(\widehat C).
 \end{aligned}
\]
回忆 $E_r^{p,q}$ 自然是 $\Lambda_{(\omega,L)}^{(0)}$-模。乘以
$e^{\pm1}\in\Lambda_{(\omega,L)}$ 定义映射
\[
 e^{\pm1}:E_r^{p,q}(\widehat C)\to E_r^{p\pm2,q}(\widehat C),
\]
从而使
$E_r(\widehat C):=\bigoplus_{p,q}E_r^{p,q}(\widehat C)$ 成为
$\operatorname{gr}_*(F\Lambda_{(\omega,L)})$-模，并且
\[
 \bigoplus_{p,q}\delta_r^{p,q}:
 \bigoplus_{p,q}E_r^{p,q}(\widehat C)
 \longrightarrow
 \bigoplus_{p,q}E_r^{p+1,q+r-1}(\widehat C)
\]
是 $\operatorname{gr}(F^*\Lambda_{(\omega,L)})$-模同态。

\begin{lemma}
\label{lem:spectral-sequence-differential}
设 $E_r^{p,q}(\widehat C)$ 是上面定义的 $\Lambda_{(\omega,L)}^{(0)}$-模。存在
$\Lambda_{(\omega,L)}^{(0)}$-模同态
\[
 \delta_r^{p,q}:E_r^{p,q}(\widehat C)\to
 E_r^{p+1,q+r-1}(\widehat C),
\]
满足
\begin{enumerate}[label=(\roman*)]
\item $\delta_r^{p+1,q+r-1}\circ\delta_r^{p,q}=0$；
\item $\ker(\delta_r^{p,q})/\operatorname{im}(\delta_r^{p-1,q-r+1})
 \cong E_{r+1}^{p,q}(\widehat C)$；
\item $e^{\pm1}\circ\delta_r^{p,q}
 =\delta_r^{p\pm2,q}\circ e^{\pm1}$。
\end{enumerate}
\end{lemma}

\begin{proof}
定义 $\delta_r^{p,q}[x]=[\delta x]\in E_r^{p+1,q+r-1}(\widehat C)$，结论立即成立。
\end{proof}

下一个引理反映了我们工作在单调情形这一事实。

\begin{lemma}
\label{lem:spectral-sequence-collapse}
谱序列在 $E_{N_L}$ 页坍缩，其中
\[
 N_L=\left[\frac{n+1}{\Sigma_L}\right].
\]
特别地，谱序列至多经过 $N_L$ 步就稳定。
\footnote{译者注：原书在本引理末句、随后关于 $E_\infty$ 的等式以及
定理 17.3.7 中均印作未定义的 $N$；本引理所定义并使用的是 $N_L$，
译文统一据此改正。}
\end{lemma}

\begin{proof}
这是~\eqref{eq:17-delta-prime-decomposition} 中给出的 $\delta$ 的结构的直接推论。
\end{proof}

这个引理使我们可以证明
\[
 E_\infty^{p,q}(\widehat C)=E_{N_L}^{p,q}(\widehat C).
\]

\begin{remark}
\label{rem:nonmonotone-projective-limit}
(FOOO09) 的一般性超出单调情形；对应谱序列并不在有限页坍缩，但其中证明了投射极限
\[
 \varprojlim E_r^{p,q}(\widehat C)
\]
存在，因而谱序列仍收敛。一般情形下该收敛性的证明见 (FOOO09)。
\end{remark}

现在令 $(\overline C,\overline\delta)$ 为 $(L,\phi_{\epsilon f}^1(L))$ 的局部\personFloer{}上同调复形；按定义，它同构于\personMorse{}复形 $(C^*(f),\delta_{\epsilon f})$。于是：

\begin{lemma}
\label{lem:E2-local-floer-morse}
存在同构
\[
 E_2^{*,*}(\widehat C)\cong
 H(\overline C;\overline\delta)\otimes_R
 \operatorname{gr}_*(F\Lambda_{(\omega,L)})
\]
作为 $\operatorname{gr}_*(F\Lambda_{(\omega,L)})$-模。
\end{lemma}

\begin{proof}
按定义，
\[
 E_1^{*,*}(\widehat C)\cong
 \overline C\otimes_R\operatorname{gr}_*(F\Lambda_{(\omega,L)}).
\]
由命题 17.1.3，$\delta_{(0)}=\overline\delta$。证明完毕。
\footnote{译者注：原书此处印作 $\delta_{(1)}=\overline\delta$；但局部\personFloer{}微分
此前记作 $\delta_{(0)}$，译文据此改正下标。}
\end{proof}

\begin{definition}
\label{def:filtration-on-homology}
定义 $F^qH(\widehat C,\delta)$ 为
$H(F^q\widehat C,\delta)$ 在 $H(\widehat C,\delta)$ 中的像。
\end{definition}

于是，由谱序列的一般结果，或由当前在有限页稳定的情形中的显式识别，得到：

\begin{theorem}
\label{thm:finite-page-associated-graded}
\[
 E_{N_L}^{p,q}(\widehat C)\cong
 F^qH^p(\widehat C,\delta)/F^{q+1}H^p(\widehat C,\delta)
\]
\footnote{译者注：原书分母漏印上标 $p$，因而不是分子所在分次部分的子模；
译文按同一定理的分子及定理 17.3.1(2) 补上 $p$。}
作为
$\Lambda^{(0)}(\lambda_0):=\Lambda_{(\omega,L)}^{(0)}
/F^{\lambda_0}\Lambda_{(\omega,L)}^{(0)}$-模。
\end{theorem}

按定义，
\[
 H^p(\widehat C,\delta)=
 HF^p(L,\phi_{\epsilon f}^1(L);\Lambda_{(\omega,L)}).
\]

\begin{remark}
\label{rem:novikov-residue-field}
回忆，在单调情形中 $\Lambda_{0,(\omega,L)}$ 是离散赋值环，其剩余域为
$\Lambda_{(\omega,L)}$。因此有同构
\begin{equation}
 \begin{aligned}
 HF(L,\phi_{\epsilon f}^1(L);\Lambda_{(\omega,L)})
 &\cong HF(L,\phi_{\epsilon f}^1(L);\Lambda_{0,(\omega,L)})\\
 &\qquad\otimes_{\Lambda_{0,(\omega,L)}}\Lambda_{(\omega,L)}.
 \end{aligned}
\label{eq:17-novikov-base-change}
\end{equation}
对子 $(L,\phi_{\epsilon f}^1(L))$，取子环 $\Lambda_{0,(\omega,L)}$ 上的\personFloer{}复形；即使在非单调情形中，也总能像对完整\personNovikov{}环
$\Lambda_{(\omega,L)}$ 那样构造它。但相伴同调
$HF(L,L';\Lambda_{0,(\omega,L)})$ 未必在哈密顿同痕下不变。另一方面，我们已经
证明，在以下情形中，单调对子 $(L,L')$ 的
$HF(L,L';\Lambda_{(\omega,L)})$ 在哈密顿同痕下不变：
\begin{enumerate}[label=(\arabic*)]
\item $L$ 与 $L'$ 的最小马斯洛夫数都大于 2（定理 16.4.6）；或
\item $L'$ 是 $L$ 的哈密顿形变，且 $L$ 的最小马斯洛夫数大于或等于 2
（定理 16.4.10）。
\end{enumerate}
\end{remark}

考察复形 $CF(L,\phi_{\epsilon f}^1(L))$ 的 $\delta$ 及这里构造的相伴谱序列，得到
作者在 (Oh96b) 中证明的如下定理。定理陈述中，在系数环 $\mathbb Z_2$ 上，我们把
$HF^*(L,\phi_{\epsilon f}^1(L))$ 简记为 $HF^*(L,L)$ 或
$HF(L;\Lambda_{(\omega,L)}^{\mathbb Z_2})$。

\begin{theorem}[与 (Oh96b) 比较]
\label{thm:large-minimal-maslov-HF}
设 $L\subset(M,\omega)$ 单调，且 $\Sigma_L\geq n+1$。考虑 $\mathbb Z_2$ 上的\personNovikov{}环 $\Lambda_{(\omega,L)}^{\mathbb Z_2}$ 及相伴\personFloer{}上同调
$HF^*(L,L)$。则：
\begin{enumerate}[label=(\arabic*)]
\item 若 $\Sigma_L>n+1$，则
$HF^*(L,L)\cong H^*(L;\mathbb Z_2)\otimes
\Lambda_{(\omega,L)}^{\mathbb Z_2}$；
\item 若 $\Sigma_L=n+1$，则对 $1<\ell<n$，
$HF^\ell(L,L)\cong H^\ell(L;\mathbb Z_2)$，并且或者
$HF^0(L,L)=HF^n(L,L)=\{0\}$，或者
\begin{equation}
 HF^n(L,L)=H^n(L;\mathbb Z_2),
 \qquad HF^0(L,L)=H^0(L;\mathbb Z_2).
\label{eq:17-end-degree-alternative}
\end{equation}
\end{enumerate}
\end{theorem}

\begin{proof}
因为 $n\geq1$，有 $\Sigma_L\geq2$，故 $HF(L,\phi(L))$ 良定义，并在 $\phi$ 的
哈密顿同痕下不变。以下简记为 $HF(L,L)$。

若 $\Sigma_L>n+1$，则 $N_L=0$，即 $\delta=\delta_{(0)}$。引理 17.3.3 与定理
17.3.7 蕴含
\[
 \begin{aligned}
 HF^*(L,L;\Lambda_{(\omega,L)}^{\mathbb Z_2})
 &\cong E_2^{*,*}(\widehat C)\\
 &\cong H(C(f);\delta_f)\otimes_{\mathbb Z_2}
 \operatorname{gr}_*(F\Lambda_{(\omega,L)}^{\mathbb Z_2})\\
 &\cong H^*(L;\mathbb Z_2)\otimes_{\mathbb Z_2}
 \Lambda_{(\omega,L)}^{\mathbb Z_2}
 \end{aligned}
\]
作为 $\Lambda_{(\omega,L)}^{\mathbb Z_2}$-模，故 (1) 得证。

若 $\Sigma_L=n+1$，则 $N_L=1$，故
$\delta=\delta_{(0)}+\delta_{(1)}\otimes q$，其中
$\delta_{(1)}:C^*(f)\to C^{*+n}(f)$。所以唯一非平凡贡献发生在 $C^0(f)$ 上。
在 $E_2$ 页，它诱导上边界映射
\begin{equation}
 \delta_{(1)}\otimes\Lambda_{(\omega,L)}^{(0);\mathbb Z_2}:
 H^0(f;\mathbb Z_2)\otimes\Lambda_{(\omega,L)}^{(0);\mathbb Z_2}
 \longrightarrow
 H^n(f;\mathbb Z_2)\otimes\Lambda_{(\omega,L)}^{(0);\mathbb Z_2},
\label{eq:17-E2-end-degree-map}
\end{equation}
并在其他所有次数上为零。因此 (2) 的第一个断言成立。对第二个断言，分
$\delta_{(1)}$ 是否为零两种情形。若不为零，它诱导同构，故
\[
 HF^0(L,L)=HF^n(L,L)=\{0\}.
\]
若为零，则~\eqref{eq:17-end-degree-alternative} 成立。对于（相对）spin 的 $L$，
上述定理可在 $\mathbb Z$ 系数上加强。
\end{proof}

最后简要提及 (FOOO09) 定义的\personFloer{}上同调上的乘积结构
\[
 \mathfrak m_2:HF(L,L;\Lambda_{(\omega,L)})\otimes
 HF(L,L;\Lambda_{(\omega,L)})\to HF(L,L;\Lambda_{(\omega,L)}).
\]
当 $\epsilon>0$ 充分小时，\personFloer{}上同调的\personMorse{}--\personBott{}版本记为
$HF(L,L;\Lambda_{(\omega,L)})$；可将它视为 $HF(L,\phi_{\epsilon f}^1(L))$。于是第
16.5 节构造的三角形乘积
\[
 \mathfrak m_2:HF(L,\phi_{\epsilon f}^1(L))\times
 HF(\phi_{\epsilon f}^1(L),L)\times HF(L,L)
\]
在令 $\epsilon\to0$ 后给出乘积
$\mathfrak m_2:\widehat C(L)\otimes\widehat C(L)\to\widehat C(L)$。由构造容易看出
\[
 \mathfrak m_2(F^{q_1}\widehat C(L)\otimes F^{q_2}\widehat C(L))
 \subseteq F^{q_1+q_2}\widehat C(L).
\]
这说明 $HF(L,L;\Lambda_{(\omega,L)})$ 上上述滤过由 $\mathfrak m_2$ 保持。下述
定理在 (FOOO09) 第 6.4.5 节中证明，证明见该处（另见 (Bu10)）。

\begin{theorem}
\label{thm:spectral-sequence-ring-compatible}
定理 17.3.1 构造的谱序列与环结构相容，具体如下。谱序列的每一页 $E_r$ 都有环结构
$\mathfrak m_2$，满足
\[
 \delta_r(\mathfrak m_2(x,y))
 =-\mathfrak m_2(\delta_r(x),y)
 +(-1)^{\deg x}\mathfrak m_2(x,\delta_r(y)).
\]
滤过 $F$ 与 $\mathfrak m_2$ 诱导的环结构相容。定理 17.3.1 中的同构都是环同构。
\end{theorem}

关于这个谱序列的进一步加强及其他应用，另见 (BCo09)、(D09)。
\stfhIndexRangeEnd{I179}{Puxulie}{谱序列（spectral sequence）}
